\documentclass[a4paper,11pt,fontset=fandol]{ctexart}
\usepackage[margin=25mm]{geometry}
\usepackage{amsmath,graphicx}
\pagestyle{plain}
\input{assets/numbers.tex}
\begin{document}
\begin{center}{\Large\bfseries 静态区间求和：已有结果的说明}\end{center}
\section{问题与接口}
本作业处理固定数组上的区间求和。
\verb|query(left, right)| 返回左闭右开区间 $[left,right)$ 的元素和。
输入为 $[\InputValues]$；本文沿用随附记录中的输出，不重新执行程序或增加计时实验。

\section{方法与边界}
\label{sec:method}
首先构建前缀数组，其次处理区间查询。令数组长度为 $n$，则
\begin{equation}
P[0]=0,\quad P[k+1]=P[k]+a[k]\quad (0\leq k<n).
\label{eq:prefix}
\end{equation}
查询公式为
\begin{equation}S(left,right)=P[right]-P[left].\label{eq:query}\end{equation}
程序检查 $0\leq left\leq right\leq n$，否则抛出 \texttt{ValueError}。
当 $left=right$ 时，两个相同的前缀值相减，空区间和为零。
下面是需要保留的接口检查代码，示例文本不是一次新的执行记录。
\begin{verbatim}
if not 0 <= left <= right <= n:
    raise ValueError("invalid range")
return prefix[right] - prefix[left]
\end{verbatim}
预处理时间和额外存储均为 $O(n)$；查询在通常整数运算模型下为 $O(1)$。
这是代码分析，不是实测结论；Python 大整数的运算成本仍受数值位数影响。
只有一次查询时，额外构建前缀数组未必值得。

\clearpage
\section{样例结果与解释}
\label{sec:results}
据随附记录\cite{record:prefix}，前缀值为 $[\PrefixValues]$。
图\ref{fig:prefix}与表\ref{tab:queries}沿用同一份既有输出；该记录的生成过程没有在本次润色中独立重跑。

\begin{figure}[!ht]
\centering\includegraphics[width=.62\linewidth]{assets/prefix.pdf}
\caption{固定输入的前缀值；横轴为元素个数，纵轴为整数和。}
\label{fig:prefix}
\end{figure}
\begin{table}[!ht]\centering\input{assets/table.tex}
\caption{三个查询的独立答案与随附输出。}\label{tab:queries}\end{table}

$[0,4)$、$[1,3)$、$[2,2)$ 的输出分别为 $\FullSum$、$\MiddleSum$、$\EmptySum$。
三个给定查询均与独立手算一致。
这仅支持所列样例，负数、越界异常和大输入尚未由随附运行验证；不能据此断言所有输入均正确。
记录未提供以 ms 为单位的耗时，本报告不比较速度。图的递增与这组正元素对应，不代表任意输入都递增。

\section{个人思考与限制}
\label{sec:thoughts}
我觉得统一成左闭右开的区间后，空区间可以直接用相同的两个前缀值相减，
不需要另加求和分支。前缀数组会多占一个元素和一份存储；如果只有一次查询，
直接遍历区间更简单。若数组经常修改，当前静态前缀和要重新构建，
是否换成其他结构还需要根据更新与查询的次数决定。

第\ref{sec:method}节的分析以静态数组为前提。后续可能比较动态更新结构，目前尚未进行该比较。
\begin{thebibliography}{9}
\bibitem{record:prefix} 随附前缀和样例输出与图表记录（由提供者交付，本次未重跑生成过程）。
\end{thebibliography}
\end{document}
